% Draft: new ending of the Introduction (Klaus, 2026-11-24). Replaces the two paragraphs
% "We focus on weak stimulation-induced changes ..." and "The framework of
% Fig.~\ref{fig:concept} ...". Content preserved; acronyms JR and QIF defined; adds a
% "Here we" statement and ends on the receiver picture instead of the calibration caveat.

Here we develop the receiver picture of Fig.~\ref{fig:concept}: cellular mixing followed
by recurrent gain, extending models of endogenous cross-frequency coupling
\cite{ruffini2025cfc,ruffini2025ham}. We work at the population level, with weak
stimulation-induced changes around ongoing activity that may already include spiking.
Direct carrier-evoked spiking and conduction block
\cite{opancar2025,peripheral2023,semenov2025} lie outside this expansion, and the
finite-amplitude network examples are identified as such.

Weak static detection produces two sources: a beat proportional to the product of the
carrier polarizations, and a steady, direct-current-like source proportional to their
summed squares. Cassar\`a et al.\ showed that this product law fixes the spatial
pattern of the beat, so that changing the current ratio only rescales it
\cite{cassara2025part2}. What is new here concerns the fate of both sources in neural
tissue. The steady source is a stimulus in its own right: it can act at depth, and its
sign is set by the state and dynamics of the detector rather than by carrier polarity.
In the exact mean field of quadratic integrate-and-fire (QIF) neurons, fast carriers
rectify, so the steady source follows the slope of the rate curve even where its
curvature predicts the opposite sign, and an unmodulated carrier matched in
root-mean-square input shows that this shift alone can drive the circuit into gamma
oscillation without any beat. Detection and amplification are separable: in a
Jansen--Rit (JR) neural mass column held at a fixed operating point, recurrent gain
varies independently of detector curvature, grows near a Hopf bifurcation as it does for
transcranial alternating-current stimulation, and appears mainly as beat-locked
modulation of the population rate, with little change in mean rate. In a realistic head
model, the current ratio leaves the beat pattern unchanged, as the product law requires,
but moves the envelope and mean-square foci in opposite directions, so where TI acts
depends on which feature the tissue detects. The size of these effects in vivo depends
on the field-to-detector transfer at kHz frequencies, which remains to be calibrated.
Within that limit, the cortical column behaves as a tuned receiver whose response is set
in turn by access, detection, and gain.
